\documentclass{article}
\usepackage[top=1.5in, bottom=1in, left=1.25in, right=1.25in]{geometry}

\usepackage{setspace}

\title{\textbf{Observability-Aware Trust-Based Grayhole Detection for Underwater Acoustic Networks}}
\date{}

\begin{document}

\setstretch{1.5}

\maketitle

\section*{Team Members}

\noindent 1. 241CS110, Appaji Nagaraja Dheeraj, appajidheeraj.241cs110@nitk.edu.in, 9481660290 \\
\noindent 2. 241CS111, Aryan Bokolia, aryanbokolia.241cs111@nitk.edu.in, 7982212773 \\
\noindent 3. 241CS117, Karthikeya Gupta, karthikeyagupta.241cs117@nitk.edu.in, 6303223736

\section*{Abstract}

% Background (4-5 lines)
Radio waves fade within a few metres in seawater, so underwater sensor networks send data with acoustic modems.
Sound travels at about 1500~m/s, links run at a few kbps, and noise, multipath and collisions corrupt many frames~\cite{akyildiz2005}.
A seabed sensor rarely reaches the surface sink directly, so other sensors relay its packets over several hops.
This routing happens at the network layer, and every node depends on its neighbours to forward honestly.
% Challenges and issues (4-5 lines)
A captured relay can launch a grayhole attack, forwarding some packets and dropping the rest.
Encryption does not stop it, because the node still holds valid keys.
Honest acoustic links already lose packets, so a relay that drops a quarter of its traffic looks like a weak link.
The usual check is a watchdog, in which a node overhears whether its next hop forwards each packet and flags the relay once its missed forwards cross a fixed threshold~\cite{marti2000}.
Underwater, the watchdog often cannot hear, because its half-duplex modem is deaf while it transmits and collisions or corrupted frames hide real forwards.
% Existing approaches and their issues (4-5 lines)
Earlier work in UnetStack3 demonstrated a selective forwarding attack and a way to detect it [3].~\cite{chandavarkar2021}.
Channel-aware detection sets a per-link threshold from expected loss, but it targets terrestrial mesh networks and does not consider times when the monitor cannot listen~\cite{shila2009}.
Underwater trust models such as ARTMM~\cite{han2015} and the model of Liu et al.~\cite{liu2024} combine several trust measures to cope with unreliable acoustic links.
In the work we reviewed, no detector checks whether the monitor could hear each expected forward.
% Problem statement (2 lines)
Watchdog detectors in underwater acoustic networks treat every missed overhear as a malicious drop, even when the monitor was transmitting or its reception failed.
These false accusations grow with traffic and let a grayhole relay hide inside natural link loss.
% Objectives (2-3)
First, we will build a watchdog in UnetStack~\cite{chitre2014} that counts a missed forward against a relay only when the monitor was able to hear it.
Second, we will compare each relay's missed forwards with the loss expected on its link, weight this evidence by link distance, and give an administrator a ranked list of suspected relays.
Third, we will compare the detector with a fixed-threshold watchdog and a channel-aware detector, measuring false accusations, detection rate, detection delay and packet delivery ratio as traffic load and drop rate vary.

\begin{thebibliography}{7}
\raggedright

\bibitem{akyildiz2005} I. F. Akyildiz, D. Pompili, and T. Melodia, ``Underwater acoustic sensor networks: research challenges,'' \emph{Ad Hoc Networks}, vol.~3, no.~3, pp.~257--279, 2005, doi: 10.1016/j.adhoc.2005.01.004.

\bibitem{marti2000} S. Marti, T. J. Giuli, K. Lai, and M. Baker, ``Mitigating routing misbehavior in mobile ad hoc networks,'' in \emph{Proc. 6th ACM MobiCom}, 2000, pp.~255--265, doi: 10.1145/345910.345955.

\bibitem{chandavarkar2021} B. R. Chandavarkar and N. Suresh, ``Demonstration of selective forwarding attack and detection mechanism in UnetStack3,'' in \emph{Proc. 12th Int. Conf. Computing Communication and Networking Technologies (ICCCNT)}, 2021, pp.~1--6, doi: 10.1109/ICCCNT51525.2021.9580125.

\bibitem{shila2009} D. M. Shila, Y. Cheng, and T. Anjali, ``Channel-aware detection of gray hole attacks in wireless mesh networks,'' in \emph{Proc. IEEE GLOBECOM}, 2009.

\bibitem{han2015} G. Han, J. Jiang, L. Shu, and M. Guizani, ``An attack-resistant trust model based on multidimensional trust metrics in underwater acoustic sensor network,'' \emph{IEEE Trans. Mobile Computing}, vol.~14, no.~12, pp.~2447--2459, 2015, doi: 10.1109/TMC.2015.2402120.

\bibitem{liu2024} C. Liu, J. Ye, F. An, and W. Jiang, ``An adaptive trust evaluation model for detecting abnormal nodes in underwater acoustic sensor networks,'' \emph{Sensors}, vol.~24, no.~9, art.~2880, 2024, doi: 10.3390/s24092880.

\bibitem{chitre2014} M. Chitre, R. Bhatnagar, and W.-S. Soh, ``UnetStack: an agent-based software stack and simulator for underwater networks,'' in \emph{Proc. IEEE OCEANS}, 2014, doi: 10.1109/OCEANS.2014.7003044.

\end{thebibliography}

\end{document}
